\documentclass[10pt,a4paper]{amsart}
\usepackage[margin=19mm]{geometry}
\usepackage{amsmath,amssymb,mathtools}
\usepackage[scr=rsfs,cal=euler]{mathalfa}
\usepackage{xfrac}
\usepackage[unicode,hidelinks,bookmarks=false]{hyperref}
\newcommand{\ie}{i.e.\ }
\newcommand{\cf}{cf.\ }
\newcommand{\resp}{respectively\ }
\newcommand{\Subsection}{\subsection}
\providecommand{\nobreakdash}{\nobreak\hbox{-}}
\setlength{\emergencystretch}{2em}
\multlinegap0pt
\usepackage{mathtools,mathrsfs}
\usepackage{booktabs}
\numberwithin{equation}{section}
\newtheorem{theorem}{Theorem}[section]
\newtheorem{proposition}[theorem]{Proposition}
\newtheorem{lemma}[theorem]{Lemma}
\newtheorem{corollary}[theorem]{Corollary}
\newtheorem{definition}[theorem]{Definition}
\newtheorem{remark}[theorem]{Remark}
\newcommand{\Ccal}{\mathcal C}
\newcommand{\Gcal}{\mathcal G}
\newcommand{\D}{\mathrm d}
\newcommand{\Rstar}{R_*}
\newcommand{\Fcal}{\mathcal F}
\begin{document}
\title[Second-Level Concavity of the Riemann Xi Kernel]{Second-Level Concavity of the Riemann $\Xi$ Kernel}
\author{Michel Planat Patrick Solé [2020] Primary 30D15; Secondary 11M26, 26A51, 26D10, 11F27, 37C10 ęywords Riemann $ $-function, Jacobi theta functions, Laguerre inequalities, Turán inequalities, log-concavity, computer-assisted proof, interval arithmetic e̱gin document; Michel Planat; Patrick Solé}
\date{}
\maketitle
\noindent\emph{Source and licence.} Second-Level Concavity of the Riemann $\Xi$ Kernel. Version arXiv:2608.19160v1 (CC BY 4.0). This material is distributed under the Creative Commons Attribution 4.0 International licence, http://creativecommons.org/licenses/by/4.0/, as recorded in the arXiv OAI licence metadata for this exact version and shown on the abstract page https://arxiv.org/abs/2608.19160v1. Full rights notice: licenses/arxiv-2608.19160-CC-BY-4.0.txt.\par\smallskip
\noindent\textit{Editorial scope.} Counted unit: the original \texttt{theorem} environment \texttt{thm:Rsharp} at source lines 433--439 together with its complete proof at lines 440--443; the delivered document keeps source lines 290--444 so that every referenced label and every citation resolves inside it. Native editable source is the paper's own concatenated TeX; the AMS wrapper is editorial formatting only. No publisher class, upstream script or unrelated artwork is redistributed.
\section{Proof II: Jacobi differential phase space}
Proof~II uses the special differential geometry of the Jacobi theta constant rather than direct control of $\Gcal$.

\subsection{Theta operator and Jacobi's third-order equation}
Let
\begin{equation}\label{eq:theta3}
 \theta_3(x)=\sum_{n\in\mathbb Z}e^{-\pi n^2x},\qquad x>0.
\end{equation}
We use the same real-nome convention for the companion theta constants,
\begin{equation}\label{eq:theta24}
 \theta_2(x)=\sum_{n\in\mathbb Z}e^{-\pi(n+1/2)^2x},\qquad
 \theta_4(x)=\sum_{n\in\mathbb Z}(-1)^n e^{-\pi n^2x}.
\end{equation}
With $x=e^{4r}$, direct differentiation gives
\begin{equation}\label{eq:thetaop}
 \boxed{\quad
 \Phi(r)=x^{5/4}\left(x\theta_3''(x)+\frac32\theta_3'(x)\right).
 \quad}
\end{equation}
Romik \cite[Eq.~(28)]{Romik2020} records Jacobi's nonlinear third-order equation for $\theta_3$.\par With $Y(x)=\theta_3(x)$, it reads
\begin{equation}\label{eq:JacobiODE}
\begin{aligned}
&\left(Y^2Y'''-15YY'Y''+30(Y')^3\right)^2
 +32\left(YY''-3(Y')^2\right)^3\\
&\hspace{20mm}=\pi^2Y^{10}\left(YY''-3(Y')^2\right)^2.
\end{aligned}
\end{equation}

Set
\begin{equation}\label{eq:LAa}
 L=\log x,\qquad \mathscr Y(L)=\theta_3(e^L),\qquad
 A=\pi e^L\mathscr Y^4,\qquad a=\frac{A'}A,
\end{equation}
and
\begin{equation}\label{eq:UVdef}
 U=2a'+1-a^2,\qquad
 V=a''+a^3-3aa'-a.
\end{equation}
Here and throughout Proof~II primes denote $L$-derivatives.

\begin{lemma}[Positivity before division]\label{lem:Upos}
Along the theta trajectory,
\begin{equation}\label{eq:Utheta}
 \boxed{\qquad U(L)=\pi^2e^{2L}\theta_2(e^L)^4\theta_4(e^L)^4>0.\qquad}
\end{equation}
\end{lemma}
\begin{proof}
With the standard modular parameter
\[
 m=\frac{\theta_2^4}{\theta_3^4},\qquad1-m=\frac{\theta_4^4}{\theta_3^4},
\]
the standard theta logarithmic-derivative identities give
$U=A^2m(1-m)$.  Since $A=\pi e^L\theta_3^4$, \eqref{eq:Utheta} follows.  This derivation does not use the variable $\chi$ introduced below.
\end{proof}

By \eqref{eq:Utheta} it is legitimate to set
\begin{equation}\label{eq:chi}
 \chi=\frac VU.
\end{equation}
Substitution of \eqref{eq:LAa}--\eqref{eq:UVdef} into \eqref{eq:JacobiODE} gives
\begin{equation}\label{eq:surface}
 4(V^2+U^3)=A^2U^2,
 \qquad
 A^2=4(U+\chi^2).
\end{equation}

\begin{definition}[Jacobi orbit]
The \emph{Jacobi orbit} is the particular curve
\[
 L\longmapsto(a(L),U(L),\chi(L))\in\mathbb R^3
\]
generated by $\mathscr Y(L)=\theta_3(e^L)$.  It is distinguished from an arbitrary solution of the autonomous system below.
\end{definition}

\begin{proposition}[Autonomous Jacobi system]\label{prop:flow}
The Jacobi orbit satisfies
\begin{equation}\label{eq:flow}
 \boxed{\begin{aligned}
 a'&=\tfrac12(U+a^2-1),\\
 U'&=2U(a+\chi),\\
 \chi'&=a\chi-U.
 \end{aligned}}
\end{equation}
Moreover
\begin{equation}\label{eq:init}
 a(0)=\chi(0)=0,\qquad
 U(0)=\Rstar:=\frac{\Gamma(1/4)^8}{64\pi^4}.
\end{equation}
\end{proposition}
\begin{proof}
The differential equations follow from \eqref{eq:UVdef}--\eqref{eq:surface}.  The classical Jacobi modular identities (see, e.g., \cite{DLMF20,ConwaySloane1999})
\[
 \theta_3(1/x)=\sqrt{x}\,\theta_3(x),\quad
 \theta_2(1/x)=\sqrt{x}\,\theta_4(x),\quad
 \theta_4(1/x)=\sqrt{x}\,\theta_2(x)
\]
show directly that $A$ and $U$ are even functions of $L$.  Hence $a=A'/A$ is odd.  Since $V$ is odd, $\chi=V/U$ is odd.  Thus $a(0)=\chi(0)=0$.  The classical value $\theta_3(1)=\Gamma(1/4)/(\sqrt2\,\pi^{3/4})$ gives \eqref{eq:init}.
\end{proof}

\subsection{Kernel factorization and elliptic monotonicity}
Equation \eqref{eq:thetaop} becomes
\begin{equation}\label{eq:kernelAU}
 \Phi(L/4)=\frac{A^{1/4}}{8\pi^{1/4}}
 \left[U+\frac32(a^2-1)\right].
\end{equation}

Let $K(m),E(m)$ be the complete elliptic integrals, put
\[
 K_c=K(1-m),\quad E_c=E(1-m),\quad q=m(1-m),
\]
and
\[
 X=E-(1-m)K,\qquad X_c=E_c-mK_c.
\]
The classical theta--elliptic parametrization gives
\begin{equation}\label{eq:ellipticbasic}
 e^L=\frac{K_c}{K},\qquad
 A=\frac{4KK_c}{\pi},\qquad
 m'=-Aq,
\end{equation}
and
\begin{equation}\label{eq:phaseell}
 1-a=\frac{4K_cX}{\pi},\qquad
 1+a=\frac{4KX_c}{\pi},\qquad
 U=A^2q.
\end{equation}
Define
\begin{equation}\label{eq:Rdef}
 R=\frac{U}{1-a^2}
 =\frac{KK_cq}{XX_c}.
\end{equation}

Put
\begin{equation}\label{eq:rhoell}
 \rho(m)=\frac{X(m)}{mK(m)}.
\end{equation}
Then $R=[\rho(m)\rho(1-m)]^{-1}$.  The Euler integral gives
\begin{equation}\label{eq:rhohyp}
 \rho(m)=\frac12
 \frac{{}_2F_1(\frac12,\frac12;2;m)}{{}_2F_1(\frac12,\frac12;1;m)}.
\end{equation}
The Markov representation of Dyachenko--Karp \cite[Example~8]{DyachenkoKarp2022}, specialized to these parameters, implies that $1-\rho$ is absolutely monotone.  Therefore $-\log\rho$ is absolutely monotone and $\log\rho$ is strictly concave on $(0,1)$.  (The integral representation has a mild logarithmic endpoint singularity after the usual real-axis continuation; for numerical evaluation the substitution $u=e^{-s}$ removes the misleading slow convergence at $u=0$.)


\begin{theorem}[Sharp elliptic ratio]\label{thm:Rsharp}
For $L\ge0$,
\begin{equation}\label{eq:Rsharp}
 \boxed{\qquad R(L)\ge\Rstar=\frac{\Gamma(1/4)^8}{64\pi^4},\qquad}
\end{equation}
with equality only at $L=0$, and $R'(L)>0$ for $L>0$.
\end{theorem}

\begin{proof}
Concavity of $g=\log\rho$ gives
$g(m)+g(1-m)\le2g(1/2)$.  Hence $R(m)\ge\rho(1/2)^{-2}$.  Legendre's relation and $K(1/2)=\Gamma(1/4)^2/(4\sqrt\pi)$ yield \eqref{eq:Rsharp}.  For $m<1/2$, $g'(m)>g'(1-m)$; since $m'(L)<0$, this gives $R'(L)>0$.
\end{proof}


\begin{thebibliography}{99}
\bibitem[ConwaySloane1999]{ConwaySloane1999} J.~H. Conway and N.~J.~A. Sloane, \emph{Sphere Packings, Lattices and Groups}, 3rd ed., Grundlehren der Mathematischen Wissenschaften, vol.~290, Springer-Verlag, New York, 1999. DOI: 10.1007/978-1-4757-6568-7.
\bibitem[DLMF20]{DLMF20} NIST Digital Library of Mathematical Functions, \emph{Chapter 20: Theta Functions}, \S20.7(viii), ``Transformations of Lattice Parameter,'' National Institute of Standards and Technology.
\bibitem[DyachenkoKarp2022]{DyachenkoKarp2022} A.~Dyachenko and D.~Karp, \emph{Integral representations of ratios of the Gauss hypergeometric functions with parameters shifted by integers}, Mathematics \textbf{10} (2022), no.~20, Article~3903, 26~pp. DOI: 10.3390/math10203903.
\bibitem[Romik2020]{Romik2020} D.~Romik, \emph{The Taylor coefficients of the Jacobi theta constant $\theta_3$}, Ramanujan J. \textbf{52} (2020), no.~2, 275--290. DOI: 10.1007/s11139-018-0109-5.
\end{thebibliography}
\end{document}
